% Module 4 section file. Included by ../module4_alfven_continuum_eigenmodes.tex.
\section{Worked examples}
\label{sec:worked_examples}

\subsection{Example 1: calculate an Alfv\'en speed and frequency}
Let
\begin{equation}
B_0=\SI{2.5}{T},
\qquad
\rho=3.3\times10^{-8}\ \mathrm{kg\,m^{-3}}.
\end{equation}
Then
\begin{align}
v_A
&=\frac{B_0}{\sqrt{\mu_0\rho}}\\
&=\frac{2.5}{\sqrt{(4\pi\times10^{-7})(3.3\times10^{-8})}}\\
&\approx1.228\times10^7\ \mathrm{m\,s^{-1}}.
\end{align}
Suppose $R_0=\SI{3}{m}$, $n=2$, $m=4$, and $\iota=0.50$. Then
\begin{equation}
n-m\iota=2-4(0.50)=0.
\end{equation}
This harmonic is at a rational surface and the reduced uncoupled frequency is zero. If instead $\iota=0.45$,
\begin{equation}
|n-m\iota|=|2-1.8|=0.2,
\end{equation}
so
\begin{align}
\omega_m&=0.2\frac{v_A}{R_0}
\approx0.2\frac{1.228\times10^7}{3}
\approx8.19\times10^5\ \mathrm{s^{-1}},\\
f_m&=\frac{\omega_m}{2\pi}\approx\SI{130}{kHz}.
\end{align}
This example shows that the dimensional scale $v_A/R_0$ is multiplied by a potentially small field-line phase factor $|n-m\iota|$.

\subsection{Example 2: locate the neighboring-harmonic crossing}
Take $n=2$, $m=3$, $\iota_0=0.35$, $\iota_1=0.65$, and $a_\iota=1$. The crossing transform is
\begin{equation}
\iota_c=\frac{2n}{2m+1}=\frac{4}{7}=0.5714286.
\end{equation}
The inverse profile gives
\begin{equation}
s_c=\frac{0.5714286-0.35}{0.65-0.35}=0.7380952.
\end{equation}
Check the two signed factors:
\begin{align}
n-m\iota_c&=2-3\left(\frac{4}{7}\right)=\frac{2}{7},\\
n-(m+1)\iota_c&=2-4\left(\frac{4}{7}\right)=-\frac{2}{7}.
\end{align}
They are equal and opposite, as required.

\subsection{Example 3: compute the local avoided-crossing gap}
Suppose the uncoupled crossing angular frequency is
\begin{equation}
\omega_c=1.50\times10^6\ \mathrm{s^{-1}}
\end{equation}
and the coupling is
\begin{equation}
C=0.20\omega_c^2.
\end{equation}
At the crossing,
\begin{align}
\Omega_-&=\omega_c\sqrt{1-0.20}=0.8944\omega_c,\\
\Omega_+&=\omega_c\sqrt{1+0.20}=1.0954\omega_c.
\end{align}
Therefore
\begin{equation}
\Delta\omega=(1.0954-0.8944)\omega_c=0.2010\omega_c.
\end{equation}
The small-coupling estimate gives
\begin{equation}
\Delta\omega\approx\frac{C}{\omega_c}=0.20\omega_c,
\end{equation}
which is accurate here to about one-half percent.

\subsection{Example 4: harmonic mixing away from and at the crossing}
Let $a-b=4C$. Then
\begin{equation}
\tan2\alpha=\frac{2C}{a-b}=\frac{1}{2},
\end{equation}
so
\begin{equation}
2\alpha\approx26.565^\circ,
\qquad
\alpha\approx13.283^\circ.
\end{equation}
The lower branch has weights approximately
\begin{equation}
\cos^2\alpha\approx0.947,
\qquad
\sin^2\alpha\approx0.053.
\end{equation}
It remains mostly one harmonic. At $a=b$, $\alpha=45^\circ$ and both weights are $0.5$. Equal mixing is therefore a localized feature of the crossing region unless coupling is very strong.

\subsection{Example 5: derive one interior finite-volume row}
Let $h=0.1$, $D_{i-1/2}=2.0\times10^9\ \mathrm{s^{-2}}$, and $D_{i+1/2}=3.0\times10^9\ \mathrm{s^{-2}}$. From Eq.~\eqref{eq:interior_stencil},
\begin{align}
K_{i,i-1}&=-\frac{2.0\times10^9}{0.1^2}=-2.0\times10^{11}\ \mathrm{s^{-2}},\\
K_{i,i}&=\frac{2.0\times10^9+3.0\times10^9}{0.1^2}=5.0\times10^{11}\ \mathrm{s^{-2}},\\
K_{i,i+1}&=-\frac{3.0\times10^9}{0.1^2}=-3.0\times10^{11}\ \mathrm{s^{-2}}.
\end{align}
The row sum is zero because an interior constant function has zero gradient. At the edge, the Dirichlet penalty makes the row sum nonzero and removes the constant null mode.

\subsection{Example 6: normalize a discrete two-harmonic mode}
Suppose $N=4$, so $h=0.25$, and an unnormalized vector is
\begin{equation}
\mathbf x=
\begin{bmatrix}
1&2&1&0&\;0&1&1&0
\end{bmatrix}^{\mathsf T}.
\end{equation}
Its Euclidean squared norm is
\begin{equation}
\mathbf x^{\mathsf T}\mathbf x=1+4+1+0+0+1+1+0=8.
\end{equation}
The grid-weighted norm is
\begin{equation}
h\mathbf x^{\mathsf T}\mathbf x=0.25(8)=2.
\end{equation}
Therefore the normalized vector is
\begin{equation}
\mathbf x_N=\frac{\mathbf x}{\sqrt2}.
\end{equation}
Using only Euclidean normalization would produce an integral norm of $h=0.25$ rather than one.

\subsection{Example 7: separate residual error from discretization error}
Assume a 128-cell calculation gives $f_{128}=\SI{271.101167}{kHz}$ with residual $10^{-12}$, while a 256-cell calculation gives $f_{256}=\SI{271.117830}{kHz}$ with the same residual. The relative difference is
\begin{equation}
\frac{|f_{256}-f_{128}|}{f_{256}}
\approx6.15\times10^{-5}.
\end{equation}
The eigensolver has solved each discrete matrix extremely accurately, but the spatial discretization still shifts the frequency by tens of parts per million. A smaller algebraic residual would not remove this grid error.

\subsection{Example 8: test exact magnetic-field scaling}
Let a reference eigenfrequency be
\begin{equation}
f_1(B_0=\SI{2.5}{T})=\SI{160.504376}{kHz}.
\end{equation}
If the field is doubled to \SI{5.0}{T}, the benchmark predicts exactly
\begin{equation}
f_1=2(160.504376)=\SI{321.008752}{kHz}
\end{equation}
up to numerical roundoff, because every matrix coefficient scales as $B_0^2$.

\subsection{Example 9: determine whether a frequency criterion is sufficient}
Mode 2 has $f_2=\SI{271.118}{kHz}$, inside the local crossing-cell interval $[183.207,297.516]\ \mathrm{kHz}$. However, its peak is at $s\approx0.275$, its centroid is $\bar s\approx0.429$, and $W_{m=4}\approx0.920$. Therefore:
\begin{equation}
\text{frequency in local interval}
\not\Rightarrow
\text{strong localization at }s_c.
\end{equation}
The correct conclusion is that the mode is useful for coupled-harmonic verification, not that frequency alone proves a trapped TAE.

\subsection{Example 10: construct a PINN trial function satisfying both boundaries}
For one component, let a neural network output $N_\theta(s)$. A trial function
\begin{equation}
\widehat\xi(s)=(1-s)N_\theta(s)
\end{equation}
satisfies the edge Dirichlet condition but not necessarily the axis Neumann condition. To satisfy both exactly, use a function of $s^2$:
\begin{equation}
\boxed{\widehat\xi(s)=(1-s^2)N_\theta(s^2).}
\end{equation}
At $s=1$, $\widehat\xi=0$. At $s=0$, every term in $d\widehat\xi/ds$ contains a factor of $s$, so $\widehat\xi'(0)=0$. The construction must be applied to both harmonic components, followed by normalization and mode-selection constraints.
